\paragraph{\trStriplikeEnd。}\label{sub:striplike-ends}
我们现在开始回顾\trFloer{}理论的若干基本支柱。这些都是人们熟识的材料——本章开头
所列的参考文献即可为证——不过，对\trFukayaCategory{}的讨论而言，它仍是一道很好的
热身练习。
首先，\trPointedBoundary\trRiemann\trSurface{}虽已足以表述 TQFT，但要实际地构造
它，一种稍更刚性的结构更为合用。记 $Z^\pm=\RR^\pm\times[0;1]$ 为半无限\trStrip。设
$S$ 是一个\trPointedBoundary\trRiemann\trSurface，其无穷远点为 $\Sigma^\pm$。$S$
的一组\trStriplikeEnd，由\trProper\trHolomorphic\trEmbedding{}$\epsilon_\zeta:
Z^\pm\to S$ 组成，对每个 $\zeta\in\Sigma^\pm$ 各有一个，且满足
\begin{equation*}
\epsilon_\zeta^{-1}(\partial S)=\RR^\pm\times\{0;1\}
\qquad\text{且}\qquad
\lim_{s\to\pm\infty}\epsilon_\zeta(s,\cdot)=\zeta,
\end{equation*}
并附加一条要求：诸 $\epsilon_\zeta$ 的像两两不交。$Z^\pm$ 上的\trRiemann\trSurface%
{}结构可延拓到一点\trCompactification{}$\widehat{Z}^\pm=Z^\pm\cup\{\pm\infty\}$
之上，而每个\trStriplikeEnd{}都扩张为一个\trHolomorphic\trEmbedding{}%
$\widehat{\epsilon}_\zeta:\widehat{Z}^\pm\to S$，把 $\pm\infty$ 映到 $\zeta$。
需要这些端自有好几条理由。其一，作为该理论分析框架的 $W^{1,p}$（$p>2$）范数并不
具有\trConformal{}不变性。同等重要的是，\trFloerCohomology{}的定义中出现的%
\trInhomogeneous{}项涉及 $Z=\RR\times[0;1]$ 上的标准坐标，因此要定义那些运算%
\eqref{eq:tqft-operation}，就需要在更一般的\trRiemann\trSurface{}的诸端上具备类似
的结构。还有一点：人们可以把\cref{fig:gluing-surfaces}中的\trGluing{}过程讲得更加
精确：给定 $\zeta_1^+$、$\zeta_2^-$ 与一个\trGluing{}长度 $l>0$，令
\begin{equation}\label{eq:striplike-gluing}
\begin{aligned}
S_1' &= S_1\setminus\epsilon_{\zeta_1^+}((l;\infty)\times[0;1]),\\
S_2' &= S_2\setminus\epsilon_{\zeta_2^-}((-\infty;-l)\times[0;1])\text{，以及}\\
S &= S_1\#_l S_2=(S_1'\sqcup S_2')/\sim,
\end{aligned}
\end{equation}
其中的等同是：对 $(s,t)\in[0;l]\times[0;1]$，取
$\epsilon_{\zeta_2^-}(s-l,t)\sim\epsilon_{\zeta_1^+}(s,t)$。
